\section{Experiment 51: results}
\label{sec:exp51_results}

\subsection{The high-redshift-safe coordinates remain accurate}

\Cref{tab:representation_redshift} shows the median RMS between each measured
CoG and its reconstruction from five coordinates. The readable error remains
between $0.00367$ and $0.00415\dex$ across all five snapshots. PCA is better by
$0.00018$--$0.00059\dex$, but both representations remain far below the
halo-to-profile prediction error. The central-censored coordinates therefore
solve the unresolved-radius problem without materially degrading compression.

\begin{table}[htbp]
  \centering
  \caption{Median full-CoG reconstruction RMS in dex. Lower is better. These
  are representation floors, not held-out halo predictions.}
  \label{tab:representation_redshift}
  \begin{tabular}{@{}lrrrrr@{}}
    \toprule
    Representation & $z=0.4$ & $z=0.7$ & $z=1.0$ & $z=1.5$ & $z=2.0$ \\
    \midrule
    Readable outer-mass coordinates & 0.00369 & 0.00367 & 0.00392 & 0.00415 & 0.00408 \\
    PCA, three components & 0.00334 & 0.00339 & 0.00374 & 0.00385 & 0.00349 \\
    \bottomrule
  \end{tabular}
\end{table}

\subsection{The readable map remains competitive with PCA}

\Cref{tab:parity_redshift} gives the held-out profile CRPS. At every epoch the
readable degree-2 model is slightly worse than PCA, by $0.00018$--$0.00070\dex$.
Each difference is smaller than the measured PCA fold-to-fold standard
deviation at that epoch, $0.00117$--$0.00333\dex$. The readable map therefore
passes the predefined parity gate at all five snapshots.

\begin{table}[htbp]
  \centering
  \caption{Five-fold held-out descendant-conditioned profile CRPS in dex.
  Lower is better. $\Delta$ is readable minus PCA, so positive values favor
  PCA.}
  \label{tab:parity_redshift}
  \begin{tabular}{@{}lrrrrr@{}}
    \toprule
    Quantity & $z=0.4$ & $z=0.7$ & $z=1.0$ & $z=1.5$ & $z=2.0$ \\
    \midrule
    Readable degree-2 & 0.06253 & 0.06651 & 0.07236 & 0.08067 & 0.09354 \\
    PCA degree-2 & 0.06223 & 0.06618 & 0.07167 & 0.08050 & 0.09294 \\
    $\Delta$ & 0.00030 & 0.00033 & 0.00070 & 0.00018 & 0.00060 \\
    PCA fold standard deviation & 0.00181 & 0.00169 & 0.00117 & 0.00261 & 0.00333 \\
    \bottomrule
  \end{tabular}
\end{table}

Nonlinearity becomes increasingly important toward high redshift. Relative to
the readable linear model, the degree-2 mean lowers profile CRPS by 3.42 per
cent at $z=0.4$, 8.33 per cent at $z=0.7$, 16.40 per cent at $z=1$, 21.63 per
cent at $z=1.5$, and 21.47 per cent at $z=2$. The stable profile coordinates do
not imply a redshift-independent linear halo relation.

\begin{figure}[htbp]
  \centering
  \includegraphics[width=0.98\textwidth]{figures/exp51_redshift_summary.pdf}
  \caption{Summary of the descendant-conditioned five-epoch experiment.
  Panel A tests only the CoG representation; panel B compares held-out
  probabilistic prediction; panel C illustrates coefficient evolution; panel D
  tests interpolation after removing an entire epoch.}
  \label{fig:exp51_summary}
\end{figure}

\subsection{Concurrent halo mass is a stronger high-redshift baseline}

At $z=2$, the final-descendant-mass-only model has a median
conditional-mean profile RMS of $0.18255\dex$. Concurrent $z=2$ halo mass alone lowers this to
$0.12501\dex$, an improvement of $0.05754\dex$. Complete descendant \diffmah{}
without concentration gives $0.13315\dex$, worse than concurrent mass alone by
$0.00814\dex$. Shuffling descendant MAH shape gives $0.18400\dex$, close to the
final-descendant-mass baseline. Adding $z=0.4$ concentration to descendant
\diffmah{} lowers the error to $0.13028\dex$ but does not overtake concurrent
mass. These controls warn that ``more history'' is not automatically a better
feature description when its normalization belongs to the wrong epoch.

\begin{figure}[htbp]
  \centering
  \includegraphics[width=0.98\textwidth]{figures/exp51_z2_predictions.pdf}
  \caption{Descendant-conditioned $z=2$ predictions in three concurrent
  halo-mass bins. Top panels compare median CoGs; bottom panels show the median
  fractional conditional-mean residual and its 16th--84th percentile range.
  Median trends can be accurate while individual residuals remain broad.}
  \label{fig:exp51_z2_descendant}
\end{figure}

\subsection{Descendant-conditioned coefficients pass held-epoch closure}

At $z=0.7$, $1.0$, and $1.5$, direct-fit CRPS is $0.06651$, $0.07236$, and
$0.08067\dex$. Interpolating coefficients without fitting the held epoch gives
$0.06733$, $0.07267$, and $0.08160\dex$, which is worse by only 1.23, 0.43, and
1.15 per cent. All pass the 5 per cent gate in \cref{eq:closure_gate}. Copying
the nearest available epoch gives $0.07530$, $0.08198$, and $0.10665\dex$, which
is substantially worse than interpolation.

This result establishes a continuous-redshift route for the
descendant-conditioned feature definition. It does not establish closure for
epoch-local inputs, because their normalizations and distributions change with
the fitting endpoint.

\begin{figure}[htbp]
  \centering
  \includegraphics[width=0.98\textwidth]{figures/exp51_coefficient_evolution.pdf}
  \caption{Evolution of the fitted standardized linear coefficients for all
  five readable profile coordinates. The figure is descriptive; the actual
  evidence for redshift continuity is the held-epoch predictive test, not the
  visual smoothness of these lines.}
  \label{fig:exp51_coefficients}
\end{figure}

\subsection{Epoch-local \diffmah{} is feasible and stronger at high redshift}

All 2365 galaxies receive valid \diffmah{} fits at all five truncation epochs.
The 11825 fits take 37.86 seconds in the measured run. Median mass-history
reconstruction RMS decreases from $0.0707\dex$ for fits ending at $z=0.4$ to
$0.0505\dex$ for fits ending at $z=2$. This decrease partly reflects the shorter
high-redshift fitted interval and must not be interpreted as stronger parameter
identification.

\Cref{tab:local_crps} compares the models using profile CRPS. Adding local
\diffmah{} and concurrent concentration lowers CRPS by 20.79--25.41 per cent
relative to concurrent halo mass alone. At $z=1.5$ and $z=2$, the local model
has CRPS $0.07197$ and $0.08363\dex$, compared with $0.07847$ and
$0.09107\dex$ for complete descendant \diffmah{} with the same concurrent
concentration. The epoch-local model is better by 8.28 and 8.17 per cent. The
high-redshift predictive signal therefore does not depend on access to later
halo growth.

\begin{table}[htbp]
  \centering
  \caption{Held-out profile CRPS in dex for the epoch-local analysis. Lower is
  better. The local model uses only history available by the listed epoch.}
  \label{tab:local_crps}
  \resizebox{\textwidth}{!}{%
  \begin{tabular}{@{}lrrrrr@{}}
    \toprule
    Model & $z=0.4$ & $z=0.7$ & $z=1.0$ & $z=1.5$ & $z=2.0$ \\
    \midrule
    Concurrent halo mass only & 0.08434 & 0.08266 & 0.08788 & 0.09649 & 0.11084 \\
    Local \diffmah{} & 0.07091 & 0.07186 & 0.07500 & 0.07928 & 0.09174 \\
    Local \diffmah{} + concurrent $c$ & 0.06443 & 0.06548 & 0.06726 & 0.07197 & 0.08363 \\
    Local MAH shape shuffled & 0.06806 & 0.07013 & 0.07381 & 0.08141 & 0.09881 \\
    Local concentration shuffled & 0.06798 & 0.06930 & 0.07148 & 0.07577 & 0.08670 \\
    Descendant \diffmah{} + concurrent $c$ & 0.06310 & 0.06430 & 0.06887 & 0.07847 & 0.09107 \\
    \bottomrule
  \end{tabular}}
\end{table}

The real pairing between a galaxy and its MAH shape becomes more valuable with
redshift. Relative to concurrent mass alone, real local \diffmah{} plus
concentration improves CRPS by 23.60 percentage points at $z=0.4$ and 24.55
points at $z=2$. When MAH shape is shuffled at fixed concurrent mass, the
improvements are only 19.29 and 10.85 points. The additional value of the real
pairing is therefore 4.31 percentage points at $z=0.4$ and 13.69 points at
$z=2$. The analogous real-versus-shuffled concentration contribution is 4.21
and 2.77 percentage points at those epochs.

\begin{figure}[htbp]
  \centering
  \includegraphics[width=0.98\textwidth]{figures/exp51_epoch_local.pdf}
  \caption{Epoch-local \diffmah{} analysis. Panel A reports the truncated-history
  fit error; panel B compares portable, descendant-conditioned, and control
  profile CRPS; panel C expresses the same comparisons as improvement relative
  to concurrent halo mass.}
  \label{fig:exp51_local}
\end{figure}

\begin{figure}[htbp]
  \centering
  \includegraphics[width=0.98\textwidth]{figures/exp51_epoch_local_z2_predictions.pdf}
  \caption{Direct profile validation of the portable $z=2$ model in three
  concurrent halo-mass bins. Median conditional-mean biases are generally a few
  per cent, while the broad residual bands show the much larger individual
  scatter represented by the stochastic model.}
  \label{fig:exp51_local_z2}
\end{figure}

\subsection{Concentration should be measured at the prediction epoch}

Holding descendant \diffmah{} fixed, replacing $z=0.4$ concentration with
concurrent concentration improves profile CRPS by 3.33 per cent at $z=0.7$,
4.66 per cent at $z=1$, 2.93 per cent at $z=1.5$, and 2.23 per cent at $z=2$.
Mass-conditioned shuffled concurrent concentrations are 0.10--2.30 per cent
worse than the $z=0.4$-concentration reference. The signal is therefore in the
correct halo--galaxy pairing and epoch label, not merely in adding one variable.

In a linear model, adding all concentration measurements available by the
prediction epoch improves on concurrent concentration by 3.40 per cent at
$z=0.4$ and 2.06 per cent at $z=0.7$, but by at most 0.17 per cent at $z\geq1$.
An extended concentration history is useful mainly when a substantial earlier
history is available.

\begin{figure}[htbp]
  \centering
  \includegraphics[width=0.98\textwidth]{figures/exp51_concentration_redshift.pdf}
  \caption{Concentration tests across redshift. Real concurrent concentration
  is compared with $z=0.4$ concentration and mass-conditioned shuffles. The
  final panel asks whether still earlier concentration values improve a linear
  model after concurrent concentration is known.}
  \label{fig:exp51_concentration}
\end{figure}

\subsection{The temporal residual model repairs over-coherent means}

The standardized held-out coordinate residuals yield $\rho=0.660$ in the AR(1)
approximation of \cref{eq:ar1}. For half-mass radius, the measured Spearman rank
correlation between $z=0.4$ and $z=2$ is 0.330. The conditional means have
correlation 0.591 and are therefore too persistent. The average correlation
among coherent stochastic draws is 0.228, which removes most of the excess but
undershoots the TNG value by 0.102.

For adjacent epoch pairs, TNG half-mass-radius correlations are 0.816, 0.819,
0.727, and 0.706. The coherent-draw averages are 0.787, 0.779, 0.712, and 0.630,
respectively. The single-$\rho$ model is a useful first approximation, but its
last adjacent pair and full-baseline persistence are too weak.

\begin{figure}[htbp]
  \centering
  \includegraphics[width=0.98\textwidth]{figures/exp51_cross_epoch_coherence.pdf}
  \caption{Cross-epoch residual and half-mass-radius rank correlations. The
  conditional mean is too coherent because regression suppresses individual
  scatter. AR(1)-coupled draws are much closer to TNG but slightly decorrelate
  the longest baseline too strongly.}
  \label{fig:exp51_coherence}
\end{figure}

\subsection{The symbolic corrections do not define a universal relation}

At $z=2$, the epoch-local linear model has profile CRPS $0.08804\dex$. The fixed
$z=0.4$ terms reach $0.08758\dex$, a 0.52 per cent improvement. A new nested
$z=2$ symbolic search reaches $0.08717\dex$, a 1.00 per cent improvement. The
dense degree-2 model reaches $0.08363\dex$, a 5.01 per cent improvement over
linear. Thus the new symbolic terms recover only about 20 per cent of the
available nonlinear gain.

No term recurs in four of five folds for total mass, $f_2$, or the outer radial
gap. A squared mass term recurs for $R_{20,\rm out}$; four coupled terms involving
mass, transition time, and growth index recur for the middle gap. This
target-specific structure is too complicated and too weak in decoded profile
performance to present as a compact redshift-independent law.

\begin{figure}[htbp]
  \centering
  \includegraphics[width=0.98\textwidth]{figures/exp51_symbolic_z2.pdf}
  \caption{Epoch-local symbolic residual test at $z=2$. The dense degree-2
  relation retains a clear predictive advantage over both transferred
  $z=0.4$ terms and a new stability-filtered symbolic search.}
  \label{fig:exp51_symbolic}
\end{figure}

\subsection{What \ExpFiftyOne{} established}

A readable, probabilistic conversion is viable at each tested snapshot and can
be made portable by refitting the halo assembly clock at the prediction epoch.
The high-redshift result is not an artifact of future halo information.
Descendant-conditioned coefficients pass an explicit redshift-interpolation
test, but continuous-redshift closure has not yet been demonstrated for the
epoch-local features. A simple temporal stochastic model is necessary and
approximately successful. No compact universal symbolic equation is supported.

\FloatBarrier
